\section{Related work}
Non-rigid registration has been widely studied in computer vision and image processing. The reader is referred to~\cite{TamCLLLMMSR13} for a recent survey on rigid and non-rigid registration of 3D point clouds and meshes. In the following, we focus on works that are closely related to our method.

Various optimization-based methods have been proposed for non-rigid registration. Chui et al.~\cite{Chui2000A} utilized a Thin Plate Spline (TPS) model to represent non-rigid mappings, and alternately update the correspondence and TPS parameters to find an optimized alignment. Following this approach, Brown et al.~\cite{brown2007Global} used a weighted variant of iterative closest point (ICP) to obtain sparse correspondences, and warped scans to global optimal position by TPS. Extending the classical ICP algorithm for rigid registration,  Amberg et al.~\cite{amberg2007optimal} proposed a non-rigid ICP algorithm that gradually decreases the stiffness of regularization and incrementally deforms the source model to the target model. Li et al.~\cite{li2008global} adopted an embedded deformation approach~\cite{sumner2007embedded} to express a non-rigid deformation using transformations defined on a deformation graph, and simultaneously optimized the correspondences between source and target scans, confidence weights for measuring the reliability of correspondence, and a warping field that aligns the source with the target. Later, Li et al.~\cite{li2009robust} combined point-to-point and point-to-plane metrics for the more accurate measure of correspondence.

Other methods tackle the problem from a statistical perspective.  Considering the fitting of two point clouds as a probability density estimation problem, Myronenko et al.~\cite{Andriy2010Point} proposed the Coherent Point Drift algorithm which encourages displacement vectors to point into similar directions to improve the coherence of the transformation.
Hontani et al.~\cite{hontani2012robust} incorporated a statistical shape model and a noise model into the non-rigid ICP framework, and detected outliers based on their sparsity. Jian et al.~\cite{jian2005robust} represented each point set as a mixture of Gaussians treated point set registration as a problem of aligning two mixtures. Also with a statistical framework, Wand et al.\cite{Wand2009Efficient} used a meshless deformation model to perform the pairwise alignment. Ma et al.~\cite{ma2015robust} proposed an ${L_2}E$ estimator to build more reliable sparse and dense correspondences.

Many non-rigid registration algorithms are based on $\ell_2$-norm metrics for the deviation between source and target models and the smoothness of transformation fields, which can lead to erroneous results when the ground-truth alignment induces large deviation or non-smooth transformation due to noises, partial overlaps, or articulate motions.
To address this issue, various sparsity-based methods have been proposed. Yang et al.~\cite{yang2015sparse} utilized the $\ell_1$-norm to promote regularity of the transformation. Li et al.~\cite{li2018robust} additionally introduced position sparsity to improve robustness against noises and outliers. For robust motion tracking and surface reconstruction, Guo et al.~\cite{guo2015robust} proposed an $\ell_0$-based motion regularizer to accommodate articulated motions.

Depending on the application, different regularizations for the transformation have been proposed to improve the robustness of the results. In~\cite{dyke2019non}, an as-rigid-as-possible energy was introduced to avoid shrinkage and keep local rigidity. Wu et al.~\cite{wu2019global} introduced an as-conformal-as-possible energy to avoid mesh distortion. Jiang et al.~\cite{jiang2019huber} applied a Huber-norm regularization to induce piecewise smooth transformation. 